\Askhsh{20684}{4}
\begin{ylh}{1}
    \item 3.1 Ο Κύκλος
    \item 3.2 Η Παραβολή.
\end{ylh}
\begin{wrapfigure}[20]{r}{8.5cm}
\begin{tikzpicture}
\begin{axis}[
    axis lines = middle,
    xlabel = {}, % Custom labels below
    ylabel = {}, % Custom labels below
    xmin = -2.5, xmax = 4.5,
    ymin = -4, ymax = 4,
    xtick = {-1,0,1,3},
    ytick = {},
    grid = none,
    width = 8cm, height = 7cm,
    tick label style = {font=\scriptsize},
    x axis line style={-stealth},
    y axis line style={-stealth},
]

    % Ορισμός συντεταγμένων με βάση τις ιδιότητες του προβλήματος και τις τιμές του ερωτήματος (β)
    \coordinate (O) at (0,0);
    \coordinate (E) at (1,0); % Εστία της παραβολής C: y^2 = 4x (p=1)
    \coordinate (A) at (3,{2*sqrt(3)}); % Σημείο A(xA, yA) με xA=3, yA=2*sqrt(3) (από β)
    \coordinate (B) at (-1,{2*sqrt(3)}); % Προβολή του A στη διευθετούσα x=-1
    \coordinate (G) at (-1,{-2*sqrt(3)}); % Σημείο Γ στη διευθετούσα, ώστε το E να είναι μέσο του ΑΓ

    % Παραβολή C: y^2 = 4x  <=> x = y^2/4
    \addplot[maincolor, thick, samples=100, domain=-3.5:4] ({x^2/4}, x);
    \node[right, maincolor] at (axis cs:4.2, 2.5) {$C: y^2 = 4x$};

    % Διευθετούσα (δ): x = -1
    \addplot[red, very thick, samples=2, domain=ymin:ymax] (-1, x);
    \node[left, red] at (axis cs:-1, 3.8) {$(\delta)$};

    % Σχεδίαση τμημάτων για το τρίγωνο ABE και την ευθεία A-Γ
    \draw[thirdcolor, thick] (A) -- (B);
    \draw[thirdcolor, thick] (A) -- (E);
    \draw[thirdcolor, thick] (B) -- (E);
    \draw[black, thick] (A) -- (G);

    % Σχεδίαση σημείων
    \fill (O) circle (1.5pt) node[below left] {O};
    \fill (E) circle (1.5pt) node[below right] {E};
    \fill (A) circle (1.5pt) node[above right] {A};
    \fill (B) circle (1.5pt) node[above left] {B};
    \fill (G) circle (1.5pt) node[below left] {$\Gamma$};

    % Προσαρμοσμένες ετικέτες αξόνων (x', x, y', y)
    \node[left] at (axis cs:-2.5,0) {$x'$};
    \node[right] at (axis cs:4.5,0) {$x$};
    \node[below] at (axis cs:0,-4) {$y'$};
    \node[above] at (axis cs:0,4) {$y$};

    % Σήμανση ορθής γωνίας στο B (μεταξύ AB και διευθετούσας)
    \draw[thirdcolor,fill=thirdcolor!20] ($(B)+(0.3,0)$) -- ($(B)+(0.3,-0.3)$) -- ($(B)+(0,-0.3)$);

\end{axis}
\end{tikzpicture}
\caption*{\vspace{-1em}} % Για αποφυγή επιπλέον κενού κάτω από την εικόνα.
\end{wrapfigure}
Ένα σημείο $A(x_A, y_A)$ της παραβολής $C: y^2 = 4x$ με $x_A > 0$, $y_A > 0$, έχει την εξής ιδιότητα: η ημιευθεία $AE$ τέμνει την διευθετούσα $(\delta)$ στο σημείο $\Gamma$, έτσι όμως ώστε η εστία $E$ της παραβολής $C$, να είναι το μέσο του τμήματος $A\Gamma$. Επίσης, από το σημείο $A$ φέρνουμε κάθετη στην διευθετούσα $(\delta)$ και έστω $B$ το σημείο τομής, όπως δείχνει το παρακάτω σχήμα.

\begin{alist}
\item Να αποδείξετε ότι το τρίγωνο $ABE$ είναι ισόπλευρο. \monades{7}
\item Να αποδείξετε ότι $x_A = 3$ και $y_A = 2\sqrt{3}$. \monades{10}
\item Να βρείτε την εξίσωση του κύκλου που διέρχεται από τις κορυφές του τριγώνου $AB\Gamma$. \monades{8}
\end{alist}